\chapter{Introduction}
\label{chap:intro}

Every organisation that has to carry out many activities with a limited set of resources faces,
sooner or later, a scheduling problem: which activity should be executed by which resource, and
when, so that deadlines are met and resources are used well. In a workshop this means deciding
the order of the operations on each machine; in a public office, which clerk processes which
file and in which order; in a project, how to arrange the tasks that share the same team. The
problem is easy to state, but hard to solve well: the number of possible plans grows
exponentially with the number of activities, and the data on which the decision is based
(durations, setup times, availability) is uncertain and changes every day.

This thesis presents \emph{FlowSync}, a web platform that supports the planner in this task. The
planner describes the problem, possibly in natural language; FlowSync turns the description into
a structured model, proposes one or more candidate plans together with their indicators, and
leaves the final decision to the human. The plans are produced by a machine learning model, which
is trained offline and then answers in a short time, instead of an optimisation solver running
at every request. This chapter describes the application domain studied
(Section~\ref{sec:domain}), the motivations of the work (Section~\ref{sec:motivations}), its
goals (Section~\ref{sec:goals}) and how the work was organised (Section~\ref{sec:contributions}).

\section{Application domain}
\label{sec:domain}

FlowSync addresses \emph{scheduling} in a broad sense. Instead of targeting a single industry, it
works on a model that is common to several domains, which are then obtained by changing the
meaning of the entities and the default weights of the objectives.

\subsection{A common model}
\label{sec:common-model}

The model is built on a small number of concepts:

\begin{itemize}
  \item an \textbf{activity} is a unit of work with a duration; activities belong to a
        \textbf{work order} (a job, a file, a project);
  \item a \textbf{resource} is whatever executes an activity: a machine, a person, a team, a room.
        Each activity can be executed only by a subset of the resources, the \emph{eligible}
        ones, for example those that have the required skill;
  \item a \textbf{precedence} constrains two activities of the same work order: the second cannot
        start before the first has finished;
  \item a \textbf{setup time} is the time needed to prepare a resource when it switches from one
        activity to another; it depends on the pair of activities involved;
  \item a \textbf{calendar} defines the working shifts and holidays of a resource, while an
        \textbf{unavailability} represents a maintenance, a leave or a breakdown;
  \item a \textbf{deadline} and a \textbf{priority} are attached to each work order.
\end{itemize}

Activities and precedences form a graph; the planner's problem is to choose, for each activity,
a resource and a start time such that all constraints hold. Among the feasible plans, the
quality is measured by one or more objectives: the \emph{makespan} (the time at which the last
activity ends), the \emph{weighted tardiness} (how much and how important the late orders are),
the \emph{overtime} and the total \emph{setup time}. Different domains weigh these objectives in
different ways, and FlowSync stores a \emph{domain profile} with the default weights, which the
planner can adjust.

\subsection{Manufacturing}
\label{sec:domain-factory}

The first domain is the production workshop. A work order is a production order made of a
sequence of operations (for example cutting, milling, finishing), each of which can be executed
on one of several machines of the right type. This is the \emph{flexible job shop} problem, one of
the most studied scheduling problems~\cite{brandimarte1993}. Setup times are relevant: changing
tool or material between two consecutive operations can take longer than the operation itself, and
the cost depends on the order in which operations are sequenced. Machines are not always
available, because of planned maintenance, and orders arrive continuously, so that the plan has to
be recomputed often. The demonstration workshop used in the prototype, with its CNC machines and a
handful of production orders, follows this setting.

\subsection{Public administration}
\label{sec:domain-pa}

The second domain is the public administration. Here the work orders are administrative files
that go through a sequence of steps (registration, check, approval), the resources are the
employees of the offices, and the skills define who is allowed to carry out a step. Deadlines are
often fixed by law, so that tardiness has a concrete cost, while setup times are usually absent or
negligible. The structure is the same as in the factory: only the vocabulary changes. This is
the main reason for choosing a generic model: the same engine can support domains that, at first
sight, have little in common.

\subsection{Projects}
\label{sec:domain-projects}

The third domain is project management, in which the activities of a project are linked by
precedences and compete for shared resources. This is the \emph{resource-constrained project
scheduling problem} (RCPSP), for which well-known benchmark libraries
exist~\cite{kolisch1997}. A project is a work order whose activities form a general graph rather
than a simple chain, and this makes it a useful test of the generality of the model.

\subsection{Scope}
\label{sec:scope}

The work focuses on the \emph{assignment and sequencing} of activities. Some related problems are
deliberately left out: in particular, the construction of staff rotas and the management of leave
requests are not addressed, and leave only enters the model as an unavailability of a resource.
The system also does not control the machines or collect data automatically from them: the actual
start and end of each activity are recorded by the operators through the application.
% TODO: complete with the other limits of the prototype (e.g. no integration with ERP systems).

\section{Motivations}
\label{sec:motivations}

In many small and medium organisations, planning is still done by hand, often on a spreadsheet.
This works while the number of activities is small, but it does not scale: with dozens of orders
and several machines, the planner can no longer evaluate the alternatives, and the plan is built
by experience and by the most urgent request. Each unexpected event, such as a breakdown or a
rush order, forces the planner to start again.

The classical alternatives have their own limits. \emph{Exact solvers} such as constraint
programming can find optimal or near-optimal plans, but scheduling problems are NP-hard, and the
computation time grows quickly with the size of the instance; when the plan must be recomputed
several times a day, waiting minutes or hours for each answer is not acceptable.
\emph{Dispatching rules}, such as earliest due date (EDD) or the apparent tardiness cost
rule (ATC)~\cite{vepsalainen1987}, are instantaneous and easy to explain, but they make each
decision looking only at the current state, and the quality of the resulting plan can be far from
the optimum.

Machine learning offers a middle way. A model trained offline on a large number of instances
can learn a dispatching policy that is richer than a hand-written rule, and at run time it
needs only a fast forward pass. Plans are built step by step by a simulator, so that they respect
by construction the constraints of the problem. A second use of learning concerns the data:
durations and setup times are not known exactly, and they can be estimated, with a measure of
confidence, from the history of what happened in the past.

Finally, an automatic plan is useful only if the planner trusts it. For this reason FlowSync is
designed as a \emph{decision support} tool and not as an autopilot: it shows several candidate
plans, with their indicators, and the planner approves, rejects or modifies them. Describing the
problem in natural language, and being asked for the missing information when needed, further
lowers the barrier for users who are not experts in optimisation.

The project was also carried out under a precise constraint: the optimisation engine must be a
machine learning model, and exact solvers cannot be used when the system is running. The
constraint programming solver is therefore used only offline, to produce the examples from which
the model learns.

\section{Goals}
\label{sec:goals}

The goal of the thesis is to design and implement a complete, domain-independent platform for
AI-assisted scheduling, and to evaluate it. In more detail:

\begin{enumerate}
  \item define a \textbf{generic problem model}, able to represent factory, public administration
        and project scheduling with the same concepts (Section~\ref{sec:common-model});
  \item use a \textbf{machine learning model as the optimisation engine}, trained by imitation
        of a constraint programming solver used only as an offline teacher, with the ATC and EDD
        rules as fallback and as baselines;
  \item produce plans that are \textbf{feasible by construction}, through an event-driven
        simulator that applies the decisions of the model;
  \item obtain plans whose quality is \textbf{competitive with the baselines} on benchmark
        instances, in a response time compatible with interactive use;
  \item keep the \textbf{planner in control}, with multiple candidate plans, objective indicators
        and an explicit approval step;
  \item deliver a \textbf{working web application} that integrates frontend, backend, database
        and AI service.
\end{enumerate}
% TODO: add quantitative targets once the evaluation protocol is fixed (e.g. improvement over ATC
% on weighted tardiness, maximum response time).

\section{Team organisation and contributions}
\label{sec:contributions}

FlowSync was developed by a team of three students, each responsible for one area:

\begin{itemize}
  \item \textbf{Leonardo Pierucci}: artificial intelligence and data (instance generator,
        dispatcher and baselines, training of M1 and M2, AI service);
  \item \textbf{Alessio Brutti}: backend and database (Spring Boot, MySQL, Flyway,
        integration with the AI service);
  \item \textbf{Mattia Brizi}: frontend (React, process graph editor, resource Gantt chart,
        dashboard).
\end{itemize}

The analysis of the problem, the definition of the interfaces between the three parts, the
documentation and the integration tests were carried out together.
% TODO: check with the supervisor whether the thesis requires an individual statement of
% contribution, and if so detail here what each student did.

\section{Thesis structure}
\label{sec:structure}
This chapter has described the application domain and the motivations of the work.
Chapter~\ref{chap:background} introduces the background on scheduling problems and their
solution methods, from exact solvers to machine learning. Chapter~\ref{chap:gantt} presents the
project plan, Chapter~\ref{chap:methodologies} the methodologies adopted, and
Chapter~\ref{chap:architecture} the software architecture of FlowSync.
% TODO: extend when the remaining chapters are added.
